\documentclass[10pt, a4paper]{article}

\usepackage[margin=1.5cm]{geometry}
\usepackage[utf8]{inputenc}
\usepackage[spanish,es-noshorthands]{babel}
\usepackage{amssymb, amsmath, amsbsy}
\usepackage{verbatim}
\usepackage{mathpazo}
\usepackage{tasks}
\usepackage{enumerate}
\usepackage[pdftex]{hyperref}
\hypersetup{pdftitle={Examen de Introducción a la Computación / Computación I },pdfauthor={Luis Carrillo Gutiérrez}}

\fontfamily{cmss}\selectfont
\renewcommand{\familydefault}{\sfdefault}
\pagestyle{empty}

\begin{document}
\begin{center}
\subsection*{Examen de la $\text{4}^{\text{ta}}$ Unidad de Aprendizaje / Contabilidad} 
\subsubsection*{Computación I – UPAC / 2011 - II} 
\end{center}

\noindent Nombres / Apellidos : \rule{10cm}{0.4pt} \\ 
Fecha : \rule{2cm}{0.4pt}  / \rule{2cm}{0.4pt} / \rule{2cm}{0.4pt} \hspace{3.4cm} Duración: {\small 20 minutos} \\

\begin{enumerate}
\item{Es un módulo \emph{optativo}, que genera las diversas  rutinas de daño de los virus informáticos, se activa acorde a un evento o circunstancia determinada, como el nacimiento de alguna personalidad.
\begin{tasks}(2)
\task{Módulo de defensa}
\task{Módulo de replicación}
\task{Módulo de ataque}
\task{Ninguna de las anteriores}
\end{tasks}
}
\item{¿Cuál es el número de la Ley Peruana (y su fecha de promulgación), que regula el \emph{envío de correo electrónico comercial no solicitado} o el manejo de \emph{(anti)SPAM}?
\begin{tasks}(2)
\task{Ley Nº 28493 del 11-04.2005}
\task{Ley Nº 28114 del 04-01.2009}
\task{Ley Nº 28394 del 14-06.2006}
\task{Ley Nº 28448 del 18-03.2008}
\end{tasks}
}
\item{¿Cuál \textbf{NO} es un objetivo de los virus informáticos?
\begin{tasks}(2)
\task{Alterar datos o la información en los medios de propagación}
\task{Eliminar el conocimiento existente en los medios de propagación}
\task{Propagarse y colapsar los medios de propagación e informáticos}
\task{Modificar el entorno en el cual opera un virus informático}
\end{tasks}
}
\item{Dada la siguiente situación: Una persona llama a tu oficina, alegando ser del servicio técnico en TIC que solucionará los recientes problemas de conexión, para tal efecto necesita conocer detalles como el modelo de router, el sistema operativo de las oficinas, los modelos de computadoras, etc. ¿De qué tipo de ataque informático es parte esta situación?
\begin{tasks}(2)
\task{Hoax telefónico}
\task{Spam}
\task{Ingeniería Social}
\task{Todas las anteriores}
\end{tasks}
}
\item{''Busca \emph{colapsar} los sistemas informáticos como el ancho de banda, así como auto-replicarse insensantemente``, es una propiedad de los :
\begin{tasks}(2)
\task{Virus informáticos}
\task{Caballos de troya}
\task{Gusanos (Worms)}
\task{Correos electrónicos no solicitado (SPAM)}
\end{tasks}
}
\item{Un \emph{virus informático} es un elemento que cumple el modelo DAS, es decir es:
\begin{tasks}(2)
\task{Dañino, Autoreplicable, Subrepticio}
\task{Deteriora, Ataca y Somete al usuario}
\task{Dañino, Atacante, Software}
\task{Todas las anteriores}
\end{tasks}
}
\item{Un \emph{sistema informático} comprometido por la acción de un virus informático, debe ser entendido como :
\begin{tasks}(2)
\task{Una computadora que necesita formateada}
\task{Un sistema de cómputo a ser reconfigurada}
\task{La necesidad de que el antivirus deba ser actualizado}
\task{De que los usuarios devuelvan la confianza entregada al sistema}
% \task{Todas las anteriores}
\end{tasks}
}
\item{La característica de \emph{auto-replicación} e \emph{ingreso furtivo}, provocando anomalías en las computadoras y desarrollando una función destructora, es propiedad de :
\begin{tasks}(2)
\task{Las bacterias informáticas}
\task{Un virus informático}
\task{Un caballo de troya con capacidades de infección}
\task{Ninguna de las anteriores}
\end{tasks}
}
\item{Dada la lista: RootKit(1), Hoax (2), Adware(3), Spyware (4), ordenar de mayor a menor daño al usuario. 
\begin{tasks}(2)
\task{1,2,3,4}
\task{4,3,1,2}
\task{4,3,1,2}
\task{2,3,4,1}
\end{tasks}
}
\item{Dentro de la siguiente lista de {\tt elementos nocivos} (virus informáticos/gusanos) en la historia de la informática, Identificar aquellos que conozca (o haya investigado)  y desarrolle características relevantes :
\begin{enumerate}[I.]
\item{SQL Slammer}
\item{I love you/loveletter}
\item{Jerusalem y variantes}
\item{Michelangelo}
\item{Win87.hytg32}
\item{CIH/Chernobyl}
\end{enumerate}
}
\end{enumerate}
\end{document}